\documentclass[12pt,reqno]{article}
\usepackage[usenames]{color}
\usepackage{amssymb}
\usepackage{amsmath}
\usepackage{amsthm}
\usepackage{fullpage}
\usepackage{float}
\usepackage[colorlinks=true,linkcolor=webgreen,filecolor=webbrown,citecolor=webgreen]{hyperref}
\definecolor{webgreen}{rgb}{0,.5,0}
\definecolor{webbrown}{rgb}{.6,0,0}

\newcommand{\seqnum}[1]{\href{https://oeis.org/#1}{\rm \underline{#1}}}
\newcommand{\NN}{\mathbb{N}}
\newcommand{\ZZ}{\mathbb{Z}}
\newcommand{\QQ}{\mathbb{Q}}
\newcommand{\RR}{\mathbb{R}}
\newcommand{\adm}{\mathcal{A}}
\newcommand{\rea}{\mathcal{R}}
\newcommand{\sq}[1]{\mathrm{#1}}

\begin{document}

\theoremstyle{plain}
\newtheorem{theorem}{Theorem}
\newtheorem{corollary}[theorem]{Corollary}
\newtheorem{lemma}[theorem]{Lemma}
\newtheorem{proposition}[theorem]{Proposition}

\theoremstyle{definition}
\newtheorem{definition}[theorem]{Definition}
\newtheorem{example}[theorem]{Example}
\newtheorem{conjecture}[theorem]{Conjecture}

\theoremstyle{remark}
\newtheorem{remark}[theorem]{Remark}

\begin{center}
\vskip 1cm{\LARGE\bf Prime Codes of Chess Positions \\
\vskip .1in
and the Chess Genome Sequence}
\vskip 1cm
\large
Jos\'e Angel C\'espedes Bejerano\\
\href{mailto:yulcesbeje@gmail.com}{\tt yulcesbeje@gmail.com}\\
\ \\
Dietrich Luebbert\\
\href{mailto:dilueb@web.de}{\tt dilueb@web.de}
\end{center}

\vskip .2 in

\begin{abstract}
We encode a chess position by assigning the first sixty-four primes to the squares of the
board and the integers from one to twelve to the piece types, and by taking the corresponding
product of prime powers. Listed in increasing order, the codes form a finite integer sequence
that we call the chess genome. For placements satisfying three simple necessary conditions for legality, we determine
the smallest code, which has twelve digits, and the largest, which has 614 digits. We give the exact number of such placements, a
fifty-digit integer, by number of men on the board, and we list the first
terms of the genome together with the number of terms below successive powers of ten. A
counting lemma relating promotions to captures, combined with elementary facts about checks,
turns the search for the largest code of a position reachable in a legal game into an integer
program; its solution has 543 digits and is attained by an explicit game. Two codes can
have logarithms that agree to seventy-four decimal places, so the logarithm is not a reliable
sorting key in floating point. We also compare the encoding with iterated Cantor pairing and
with base-thirteen numerals.
\end{abstract}

\section{Introduction}\label{secintro}

Shannon~\cite{shannon} estimated the number of chess positions to be roughly $10^{43}$ and
the number of possible games to be at least $10^{120}$. Tromp~\cite{tromp} estimates the number of legal positions at about
$4.8\cdot 10^{44}$. For positions without promoted men, Steinerberger~\cite{steinerberger} and
Gourion~\cite{gourion} proved upper bounds, the latter below $4\cdot 10^{37}$ diagrams.

This paper is about writing positions as integers. Number the squares from
$1$ to $64$, attach the $i$-th prime $p_i$ to square $i$, label the twelve kinds of men by
$1,2,\ldots,12$ and the empty square by $0$, and send a placement with labels
$(v_1,\ldots,v_{64})$ to the integer $\prod_i p_i^{v_i}$. G\"odel~\cite{godel} coded finite
sequences in this way; here the sequence is a word of length $64$ over an alphabet of $13$
letters, and the primes are sequence \seqnum{A000040} of the OEIS~\cite{oeis}. By unique
factorization the code is injective, and the contents of each square can be read from it as a
$p$-adic valuation.

Listing the codes in increasing order turns the set of positions into a finite increasing
sequence of integers, which we call the chess genome. We study this sequence: how long it
is, how it begins, how large its terms become, and how close together they lie.

A difficulty appears at once. Whether a given placement can occur in a game is a question of
retrograde analysis, for which no simple test is known. We therefore work with the larger set
$\adm$ of statically admissible placements, defined by three conditions that are checked
at a glance (Definition~\ref{defadm}), and we treat the reachable placements separately.

On $\adm$ the basic quantities can be computed exactly. The code $S$ and its logarithm $L$ are injective and
define the same order (Theorems~\ref{thminj} and~\ref{thmorder}). The smallest code is
$2^{12}\cdot 5^{11}=2\cdot 10^{11}$ (Theorem~\ref{thmmin}), and the largest is a $614$-digit
integer belonging to one placement with nine queens on each side (Theorem~\ref{thmmax}). The
set $\adm$ has exactly
\begin{displaymath}
N=@@NTOTAL@@
\end{displaymath}
elements, and Theorem~\ref{thmcount} gives the number $c(m)$ of them with $m$ men. Its first
value, $c(2)=3612$, is the term of index $8$ of \seqnum{A035286}. Tables~\ref{tabfirst}
and~\ref{tabcount} list the first terms of the genome and the number of terms up to $10^k$ for
$k\leq @@KMAX@@$.

For reachable placements we can also find the largest code. It has $543$ digits and belongs to
a placement with $26$ men, which an explicit game of $272$ plies reaches
(Theorem~\ref{thmreach}). The proof rests on two facts about games. The first concerns
promotions: after $a$ captures of pawns and $b$ captures of other men, at most
$\min(3a+2b,16-a)$ pawns have been promoted, and for every possible pair $(a,b)$ some game
reaches this bound (Lemma~\ref{lemcap} and Theorem~\ref{thmsharp}). The second is a short list of
constraints on checks (Lemma~\ref{lemcheck}). Together they turn the search into an integer
program, which we solve by computer. The same program bounds the code for each number of men
(Proposition~\ref{propmen}). Among the first $433$ terms of the genome, all but two belong to
reachable placements, and the two exceptions are ruled out by the constraints on checks
(Proposition~\ref{propreach}). On an $n\times n$ board, the number of placements with the two
kings and one more man is a polynomial of degree $6$ in $n$ (Proposition~\ref{propnn}).

Finally, we ask how well $L$ separates placements. The smallest difference between two values of
$L$ on $\adm$ lies between $10^{-138}$ and $1.91\cdot 10^{-75}$ (Proposition~\ref{propgap}), so a
floating-point value of $L$ cannot be trusted as a sorting key.

Section~\ref{secenc} fixes the encoding and Section~\ref{secorder} the order. Sections~\ref{secext}
and~\ref{seccap} contain the extreme values and the bound for reachable positions.
Section~\ref{secgenome} treats the genome as a sequence, Section~\ref{secother} compares the
code with two other encodings, and Section~\ref{secopen} contains concluding remarks.

\section{The encoding}\label{secenc}

\subsection{Squares, labels, and codes}

Squares are numbered row by row starting from $\sq{a1}$, so that $\sq{a1},\sq{b1},\ldots,\sq{h1}$
are squares $1,\ldots,8$, square $9$ is $\sq{a2}$, and square $64$ is $\sq{h8}$. Square $i$ carries
the prime $p_i$; thus $p_1=2$, $p_2=3$, $p_3=5$, and $p_{64}=311$. The contents of a square are
recorded by the labels of Table~\ref{tablabels}. White men have odd labels, black men even ones,
and we call rooks, knights, bishops, and queens officers.

\begin{table}[H]
\begin{center}
\begin{tabular}{l|ccccccccccccc}
contents & empty & P & p & R & r & N & n & B & b & Q & q & K & k\\ \hline
label & 0 & 1 & 2 & 3 & 4 & 5 & 6 & 7 & 8 & 9 & 10 & 11 & 12
\end{tabular}
\end{center}
\caption{Labels of the men. Capital letters denote white men and lower-case letters black men.}
\label{tablabels}
\end{table}

A placement is a vector $B=(v_1,\ldots,v_{64})\in\{0,1,\ldots,12\}^{64}$. Its code
and its logarithmic code are
\begin{equation}\label{eqcode}
S(B)=\prod_{i=1}^{64}p_i^{v_i}
\qquad\text{and}\qquad
L(B)=\ln S(B)=\sum_{i=1}^{64}v_i\ln p_i .
\end{equation}
We write placements in the form $\sq{ka1},\sq{Kc1},\sq{Pa2}$, meaning a black king on
$\sq{a1}$, a white king on $\sq{c1}$, and a white pawn on $\sq{a2}$.

\begin{remark}
Moves act on the code in a simple way. The contents of square $i$ are
the exponent of $p_i$ in $S(B)$. Moving a man with label $v$ from square $i$ to an empty
square $j$ multiplies the code by $(p_j/p_i)^v$, and a capture on square $j$ of a man with
label $u$ divides it further by $p_j^u$. In logarithmic form a move is a translation of $L$ by
$v(\ln p_j-\ln p_i)$.
\end{remark}

\subsection{Static admissibility}

\begin{definition}\label{defadm}
A placement $B$ is statically admissible if the following three conditions hold.
\begin{itemize}
\item[(S1)] Exactly one entry equals $11$ and exactly one equals $12$, and the two squares carrying
them are not adjacent, diagonal neighbors included.
\item[(S2)] The entries $1$ and $2$ occur only on squares $9,\ldots,56$, that is, pawns stand on ranks
$2$ to $7$.
\item[(S3)] For each color, if $P$ denotes the number of pawns and $R$, $N$, $B$, $Q$ the numbers
of rooks, knights, bishops, and queens of that color, then
\begin{equation}\label{eqexcess}
e=\max(R-2,0)+\max(N-2,0)+\max(B-2,0)+\max(Q-1,0)\leq 8-P .
\end{equation}
\end{itemize}
We denote by $\adm$ the set of statically admissible placements and by $\rea$ the set of
reachable placements, those occurring in some game played from the initial array
according to the laws of chess~\cite{fide}.
\end{definition}

The integer $e$ in~\eqref{eqexcess} is the excess of the color: it counts officers that
cannot belong to the original complement and must therefore be promoted pawns. A promotion
uses up a pawn, so a reachable placement satisfies (S3); it satisfies (S1) and (S2) by the
rules for kings and pawns. Hence $\rea\subseteq\adm$. The two sets are far from equal. The
conditions say nothing about checks, about the colors of the squares of the bishops, or about
pawn structure, and they ignore that promotions require captures. We come back to captures and
checks in Section~\ref{seccap}.
Condition (S3) also bounds the number of men of a color by $1+P+7+e\leq 16$.

\subsection{The complete position}

A position in the sense of the laws of chess consists of a placement together with the side to
move, the castling rights, and a possible en passant capture. We record the side to move by
$t\in\{0,1\}$, with $t=1$ if Black is to move. We record the castling rights by a pair
$(w,b)\in\{0,1,2,3\}^2$, with the value $0$ for no right, $1$ for the right to castle on the
king's side only, $2$ for the queen's side only, and $3$ for both; $w$ refers to White and
$b$ to Black. Finally $j\in\{0,1,\ldots,16\}$ records the en passant situation: $j=0$ if there
is none, $j\in\{1,\ldots,8\}$ if the white pawn of file $j$ (with $\sq{a}=1$, \ldots, $\sq{h}=8$) has
just advanced two squares, and $j\in\{9,\ldots,16\}$ if the black pawn of file $j-8$ has. The
complete code is
\begin{equation}\label{eqtotal}
S_{\mathrm{tot}}=S(B)\cdot 313^{t}\cdot \gamma(w,b)\cdot\varepsilon(j),
\end{equation}
with $\gamma(0,0)=\varepsilon(0)=1$, and otherwise
\begin{displaymath}
\gamma(w,b)=p_{65+w+4b},\qquad \varepsilon(j)=p_{80+j} .
\end{displaymath}
Thus $313=p_{65}$ marks the side to move, the fifteen primes from $p_{66}=317$ to $p_{80}=409$
mark the castling rights (Table~\ref{tabcastle}), and the sixteen primes from $p_{81}=419$ to
$p_{96}=503$ mark the en passant file. The largest possible factor is
$313\cdot 409\cdot 503=64392551<10^{8}$.

\begin{table}[H]
\begin{center}
\begin{tabular}{l|cccc}
 & $w=0$ & $w=1$ & $w=2$ & $w=3$\\ \hline
$b=0$ & $1$ & $317$ & $331$ & $337$\\
$b=1$ & $347$ & $349$ & $353$ & $359$\\
$b=2$ & $367$ & $373$ & $379$ & $383$\\
$b=3$ & $389$ & $397$ & $401$ & $409$
\end{tabular}
\end{center}
\caption{The castling factor $\gamma(w,b)$.}
\label{tabcastle}
\end{table}

\section{Injectivity and the logarithmic order}\label{secorder}

\begin{theorem}\label{thminj}
The map $B\mapsto S(B)$ is a bijection from $\{0,1,\ldots,12\}^{64}$ onto the set of positive
integers all of whose prime factors are at most $311$ and occur with exponent at most $12$.
The map $(B,t,w,b,j)\mapsto S_{\mathrm{tot}}$ is injective.
\end{theorem}

\begin{proof}
By the fundamental theorem of arithmetic~\cite[Thm.\ 2]{hardy}, a positive integer determines
the exponent of each prime in its factorization. Applied to $p_1,\ldots,p_{64}$ this gives both
assertions about $S$. In $S_{\mathrm{tot}}$ the primes $p_{65},\ldots,p_{96}$ do not divide
$S(B)$, and the three factors $313^t$, $\gamma(w,b)$, and $\varepsilon(j)$ involve disjoint sets
of primes, each value of $(w,b)$ and of $j$ corresponding to a different prime or to $1$. So
$S_{\mathrm{tot}}$ determines $S(B)$, $t$, $(w,b)$, and $j$.
\end{proof}

\begin{theorem}\label{thmorder}
For placements $A$ and $B$ we have $S(A)<S(B)$ if and only if $L(A)<L(B)$. The map $L$ is also
injective, so that $A\prec B$ defined by $L(A)<L(B)$ is a strict total order on placements.
\end{theorem}

\begin{proof}
The first assertion holds because $L=\ln S$ and the logarithm is strictly increasing. If
$L(A)=L(B)$, then $S(A)=S(B)$ and $A=B$ by Theorem~\ref{thminj}.
\end{proof}

The proof uses nothing beyond unique factorization, which also shows that the numbers
$\ln p_1,\ldots,\ln p_{64}$ are linearly independent over $\QQ$. How close two values of $L$
can be is another question, to which we return in Proposition~\ref{propgap}.

\section{The smallest and the largest code}\label{secext}

\begin{theorem}\label{thmmin}
The minimum of $S$ on $\adm$ is $2^{12}\cdot 5^{11}=200000000000$. It is attained only by the
placement $\sq{ka1},\sq{Kc1}$.
\end{theorem}

\begin{proof}
Let $B\in\adm$ have its black king on a square with prime $q$ and its white king on a square
with prime $r$. Then $S(B)\geq q^{12}r^{11}$, with equality exactly when there are no other men.
If neither king stands on $\sq{a1}$, then $q$ and $r$ are distinct primes not smaller than $3$, so
$q^{12}r^{11}\geq 3\cdot 15^{11}>2^{12}\cdot 5^{11}$. If a king stands on $\sq{a1}$, the other one is
not on $\sq{b1}$, $\sq{a2}$, or $\sq{b2}$ by (S1), so its prime is at least $5$, the prime of $\sq{c1}$.
The two possibilities give $S(B)\geq 2^{12}\cdot 5^{11}$ if the black king is on $\sq{a1}$ and
$S(B)\geq 2^{11}\cdot 5^{12}>2^{12}\cdot 5^{11}$ if the white king is.
\end{proof}

The maximum takes more work. We find the material first and the squares afterwards. The
material can be settled first because a man may be replaced, on the square where it stands, by
one with a larger label, and this raises the code whatever the rest of the board looks like.
Lemma~\ref{lemmat} shows that condition (S3) allows such replacements until a single material
is left. With the material fixed, placing the men is a sorting problem
(Lemma~\ref{lemplace}), except for the two kings: they must not be adjacent, and we simply try
all pairs of squares for them.

\begin{lemma}\label{lemmat}
If $B$ maximizes $S$ on $\adm$, then each color has a king, nine queens, two rooks, two bishops,
two knights, and no pawn.
\end{lemma}

\begin{proof}
We show that any other material can be improved inside $\adm$. Conditions (S1) and (S2) are not
affected by adding an officer on an empty square or by replacing a man by an officer on the same
square, so only (S3) needs attention.

Suppose first that some color has fewer than $16$ men. If it has fewer than two rooks, two
knights, two bishops, or no queen, we add an officer of such a kind on an empty square; the
excess $e$ does not change. Otherwise the color has $7+e$ officers, hence $8+P+e<16$ men, so
that $e+1\leq 8-P$ and we may add a queen. In both cases $S$ increases.

Suppose next that some color has a pawn. Replacing it by a queen of the same color lowers $P$
by one and raises $e$ by at most one, so~\eqref{eqexcess} persists, and the label of the square
increases.

We are left with $P=0$ and $15$ officers for each color. At most seven of them are counted
outside $e$, so $e\geq 8$, and~\eqref{eqexcess} forces $e=8$ together with $R\geq 2$, $N\geq 2$,
$B\geq 2$, and $Q\geq 1$. If $R\geq 3$, we replace a rook by a queen of the same color: $e$ is
unchanged and the label increases. The same applies to knights and bishops. Hence
$R=N=B=2$ and $Q=9$.
\end{proof}

\begin{lemma}\label{lemplace}
Fix two non-adjacent squares for the kings and, for each color, numbers of pawns and of officers
of each kind. Among the placements with these data, if there are any, $S$ is largest for the
following one: the officers, sorted by decreasing label, occupy the free squares in decreasing
order of the index, and then the pawns, black ones first, occupy in decreasing order the
remaining free squares of ranks $2$ to $7$.
\end{lemma}

\begin{proof}
Let $B$ be a maximizer. An officer cannot have an empty square above it in the numbering,
since moving it there would increase $S$. In particular no officer stands on rank $1$: otherwise
all $56$ squares of ranks $2$ to $8$ would be occupied, while there are at most $32$ men. If a
pawn stood on a square $y$ and an officer on a square $x<y$, then $x$ would lie on ranks $2$ to
$7$ and exchanging the two men would increase $S$, the label of an officer being at least $3$ and
that of a pawn at most $2$. So the officers fill the largest free squares. Exchanging two officers,
or two pawns, that are not in decreasing order of labels increases $S$ as well, and a pawn
cannot have an empty square of ranks $2$ to $7$ above it.
\end{proof}

\begin{theorem}\label{thmmax}
The maximum of $S$ on $\adm$ is attained only by the placement $B_{\max}$ with
\begin{displaymath}
\begin{array}{ll}
\sq{k}\colon \sq{h8}, & \sq{K}\colon \sq{f8},\\
\sq{q}\colon \sq{g8},\sq{e8},\sq{d8},\sq{c8},\sq{b8},\sq{a8},\sq{h7},\sq{g7},\sq{f7}, &
\sq{Q}\colon \sq{e7},\sq{d7},\sq{c7},\sq{b7},\sq{a7},\sq{h6},\sq{g6},\sq{f6},\sq{e6},\\
\sq{b}\colon \sq{d6},\sq{c6}, & \sq{B}\colon \sq{b6},\sq{a6},\\
\sq{n}\colon \sq{h5},\sq{g5}, & \sq{N}\colon \sq{f5},\sq{e5},\\
\sq{r}\colon \sq{d5},\sq{c5}, & \sq{R}\colon \sq{b5},\sq{a5} .
\end{array}
\end{displaymath}
Its code $S_{\max}=S(B_{\max})$ has $614$ decimal digits, and
\begin{displaymath}
\log_{10}S_{\max}=613.41778\ldots,\qquad L(B_{\max})=1412.44664\ldots .
\end{displaymath}
\end{theorem}

\begin{proof}
By Lemma~\ref{lemmat} the material is that of $B_{\max}$, and by Lemma~\ref{lemplace} the
placement is determined by the squares of the two kings. There are $3612$ ordered pairs of
non-adjacent squares. We evaluated the code of the corresponding placement for each of them
in exact integer arithmetic; the largest value occurs for the black king on $\sq{h8}$ and the
white king on $\sq{f8}$. Uniqueness follows from Theorem~\ref{thminj}.
\end{proof}

\begin{corollary}\label{corbounds}
For every $B\in\adm$ we have $2\cdot 10^{11}\leq S(B)<10^{614}$ and
$26.02<L(B)<1412.45$. Every complete code satisfies $S_{\mathrm{tot}}<10^{622}$.
\end{corollary}

\begin{proof}
The bounds for $S$ and $L$ restate Theorems~\ref{thmmin} and~\ref{thmmax}. For the complete
code, $\log_{10}S_{\mathrm{tot}}\leq \log_{10}S_{\max}+\log_{10}64392551<613.42+7.81$.
\end{proof}

\begin{example}
The initial array has a code with $239$ digits, $\log_{10}S=238.09547\ldots$, and
$L=548.23509\ldots$.
\end{example}

The placement $B_{\max}$ cannot occur in a game. Its white king is attacked by four queens,
whereas in a game a king is never attacked by more than two men (Lemma~\ref{lemcheck}). In the next section we find out
how large the code of a reachable placement can be.

\section{Promotions, captures, and reachable positions}\label{seccap}

\begin{lemma}\label{lemcap}
Consider a game played from the initial array, and at some moment of it let $a$ be the number
of captures made so far in which the captured man was a pawn, $b$ the number of other
captures, $\pi$ the number of promotions, and $P$ the number of pawns on the board. Then
\begin{equation}\label{eqcap}
P+\pi+a=16\qquad\text{and}\qquad \pi\leq 3a+2b .
\end{equation}
The inequality cannot be improved: there is a game with $a=1$, $b=0$, and $\pi=3$.
\end{lemma}

\begin{proof}
Each of the $16$ original pawns is still a pawn on the board, or has been promoted, or was
captured while it was a pawn. This is the equality.

For $i=1,\ldots,8$ let $W_i$ and $Z_i$ be the white and the black pawn that start on file $i$. We
say that a capture touches the index $i$ if the capturing man or the captured man is
$W_i$ or $Z_i$, still a pawn at that time. We claim that if $W_i$ or $Z_i$ has been promoted,
then some capture touches $i$. Consider the first of the two pawns to be promoted, say $W_i$,
at time $T$; the other case is symmetric. If $Z_i$ was captured before $T$, it was captured as
a pawn and we are done. Otherwise $Z_i$ is a pawn on the board at time $T$. A pawn leaves its
file only by capturing. If neither pawn had captured up to time $T$, both would have stayed
on file $i$, with $W_i$ moving up, $Z_i$ moving down, and $W_i$ initially below $Z_i$; as they
cannot pass each other, $W_i$ would still be below $Z_i$ at time $T$, which is absurd because
$W_i$ is then on the eighth rank. So one of them has made a capture as a pawn.

Now let $\pi_i\in\{0,1,2\}$ be the number of promoted pawns among $W_i$ and $Z_i$. For
each $i$ with $\pi_i>0$ choose a capture touching $i$. A capture touches at most two indices,
that of the capturing man and that of the captured man. If the captured man is a pawn with
index $j$, it is never promoted, so $\pi_j\leq 1$; together with a second index $i\neq j$ this
accounts for at most $1+2=3$ promotions. If the captured man is not a pawn, only the index
of the capturing man is touched, which accounts for at most $2$ promotions. Summing over
the chosen captures gives $\pi=\sum_i\pi_i\leq 3a+2b$.

Sharpness is shown by the game
1.\ a4 b5 2.\ axb5 Nc6 3.\ b6 a5 4.\ b7 a4 5.\ b8=Q a3 6.\ b3 Nf6 7.\ Nc3 Ng8 8.\ Rb1 a2
9.\ Rb2 a1=Q 10.\ Qb4 Nf6 11.\ Qd4 Ng8 12.\ b4 Nf6 13.\ b5 Ng8 14.\ b6 Nf6 15.\ b7 Ng8
16.\ b8=Q, in which the only capture is that of a pawn by a pawn and three pawns are
promoted.
\end{proof}

The game in the proof shows where the constants $3$ and $2$ come from. After 2.\ axb5 the
white pawn from the a-file stands on the b-file and the black b-pawn has disappeared, so three
pawns have a free road at the price of one capture: the capturing pawn, the white b-pawn behind
it, and the black a-pawn. Had the pawn taken a knight on b5 instead, the black b-pawn would
still be in the way, and only the two pawns that started on the a-file could hope to profit.

The equality in~\eqref{eqcap} also limits promotions: since $P\geq 0$, at most $16-a$ pawns are
promoted. Together with Lemma~\ref{lemcap} this gives $\pi\leq\min(3a+2b,16-a)$, and this bound is
the true maximum.

\begin{theorem}\label{thmsharp}
Let $a,b\geq 0$ be integers with $a\leq 16$ and $a+b\leq 30$. The largest number of promotions
at a moment of a game played from the initial array, after $a$ captures of pawns and $b$ other
captures, is $\min(3a+2b,16-a)$. In particular $3a+2b$ promotions occur if and only if
$2a+b\leq 8$.
\end{theorem}

The conditions on $a$ and $b$ are needed for such a moment to exist at all, since there are
$16$ pawns and $30$ men other than the kings. The first term of the minimum is the smaller one
exactly when $2a+b\leq 8$.

\begin{proof}
The upper bound was shown above, so we must exhibit games. We use local patterns, each on a
pair of adjacent files $f$ and $g$. In the first, the white pawn of file $f$ takes the black pawn of
file $g$, as in the proof of Lemma~\ref{lemcap}; then file $f$ holds only a black pawn and file $g$
only white pawns, and three pawns promote at the price of one capture of a pawn. In the
second, the black pawn of file $g$ takes a white officer on file $f$, and the white pawn of file $f$
takes a black officer on file $g$; then file $f$ holds two black pawns and file $g$ two white ones,
and all four promote at the price of two captures of officers. In the third, the black pawn of
file $g$ first advances until it stands in front of the white pawn of file $g$, and then the white
pawn of file $f$ takes a black officer on file $g$, on a square beyond the black pawn; the capturing
pawn and the black pawn of file $f$ promote, at the price of one capture of an officer. In the
last, officers take $j\in\{2,3,4\}$ of the four pawns of the two files, chosen so that each remaining
pawn has a free file; the $4-j$ remaining pawns promote.

The board has four pairs of adjacent files, namely $\sq{a}\sq{b}$, $\sq{c}\sq{d}$, $\sq{e}\sq{f}$,
and $\sq{g}\sq{h}$. Suppose first that $2a+b\leq 8$. Use the first pattern $a$ times, the second
$\lfloor b/2\rfloor$ times, and the third once if $b$ is odd, each time on a new pair. This needs
$a+\lceil b/2\rceil\leq 4$ pairs and gives $3a+4\lfloor b/2\rfloor+2(b-2\lfloor b/2\rfloor)=3a+2b$
promotions.

Suppose next that $2a+b>8$, so that the bound is $16-a$. If $a\leq 4$, play the game just
described for the pair $(a,8-2a)$; at its end every pawn that was not captured has promoted, so
there are $16-a$ promotions and no pawn is left, and we continue with $b-(8-2a)$ captures of
officers. If $a\geq 4$, write $a=j_1+j_2+j_3+j_4$ with $1\leq j_k\leq 4$ and use, on the $k$-th pair,
the first pattern if $j_k=1$ and the last one otherwise. Exactly $a$ pawns are captured and the
other $16-a$ promote; we continue with $b$ captures of officers. Captures of officers never
change $a$ or $\pi$, so one game serves all admissible values of $b$ for a given $a$.

To turn the patterns into games, the other men must clear the way: the squares in front of the
promoting pawns, the promotion squares, and the squares of the men to be taken. We did this by
a search that treats one capture or one promotion at a time and forbids all other captures. This
gave $24$ games for the pairs with $2a+b\leq 8$ and $(a,b)\neq(0,0)$, and $17$ games, one for
each $a\leq 16$, that pass through moments with $a$, $b$, and $\min(3a+2b,16-a)$ for every
admissible $b$ with $2a+b\geq 8$. The longest has $207$ plies. They are included with the
programs~\cite{software}, and each was replayed with the python-chess library~\cite{pychess},
which checks every move for legality, and recounted.
\end{proof}

The lemma is about games, but what it needs can be read off a placement: every missing man
was captured, and every missing pawn was captured or promoted. For a placement $B$ let $m$ be the number of men, $P$ the total number of pawns,
and $e$ the sum of the excesses~\eqref{eqexcess} of the two colors.

\begin{corollary}\label{corcap}
If $B\in\rea$, then there is an integer $a$ with
\begin{equation}\label{eqfeas}
0\leq a\leq 32-m,\qquad e\leq 16-P-a,\qquad 16-P-a\leq a+2(32-m) .
\end{equation}
In particular, a reachable placement with all $32$ men has no promoted man, and a reachable
placement without pawns has at most $28$ men.
\end{corollary}

\begin{proof}
A capture removes one man and nothing else does, so $a+b=32-m$ at the moment $B$
occurs. At least $e$ of the officers on the board are promoted pawns, so
$e\leq\pi=16-P-a$ by~\eqref{eqcap}, and $\pi\leq 3a+2b=a+2(32-m)$. For $m=32$ the first
condition gives $a=0$, the third then gives $P\geq 16$, and the second $e\leq 0$. For $P=0$
the third condition and $a\leq 32-m$ give $16\leq 2a+2(32-m)\leq 4(32-m)$.
\end{proof}

We now look for the largest code among the placements of $\adm$ that
satisfy~\eqref{eqfeas}. This is a finite search, and one observation keeps it small. For one
color, call $(P,O,e)$ its class, with $O$ the
number of officers. Conditions (S3) and~\eqref{eqfeas} see an army only through its class, so
inside a class we may take the officers that give the largest code. There is always a best choice,
namely $e$ promoted queens together with the $O-e$ highest of the seven original officers in the
order queen, bishops, knights, rooks: its labels, sorted decreasingly, are entry by entry at least
those of any other army of the class. We checked this over the $8694$ admissible vectors
$(P,R,N,B,Q)$, which fall into $324$ classes. Replacing an army by the dominant one of its class,
man for man on the same squares, therefore does not decrease $S$ and changes neither $m$, $P$,
nor $e$.

\begin{proposition}\label{propcap}
Let $B_{\ast}$ be the placement with $26$ men
\begin{displaymath}
\begin{array}{ll}
\sq{k}\colon \sq{h8}, & \sq{K}\colon \sq{f8},\\
\sq{q}\colon \sq{g8},\sq{e8},\sq{d8},\sq{c8},\sq{b8},\sq{a8},\sq{h7},\sq{g7},\sq{f7}, &
\sq{Q}\colon \sq{e7},\sq{d7},\sq{c7},\sq{b7},\sq{a7},\sq{h6},\sq{g6},\\
\sq{b}\colon \sq{f6},\sq{e6}, & \sq{B}\colon \sq{d6},\sq{c6},\\
\sq{n}\colon \sq{b6},\sq{a6}, & \sq{N}\colon \sq{h5},\sq{g5} .
\end{array}
\end{displaymath}
Every $B\in\adm$ satisfying~\eqref{eqfeas} for some integer $a$ has $S(B)\leq S(B_{\ast})$, and
$\log_{10}S(B_{\ast})=542.99828\ldots$.
\end{proposition}

\begin{proof}
By the preceding discussion it suffices to consider the $324^2=104976$ pairs of classes with
dominant armies, of which $102672$ satisfy~\eqref{eqfeas} for some $a$. For such a pair, let $U$
be the value of $L$ obtained by putting the labels, sorted decreasingly, on the squares
$64,63,\ldots$, ignoring (S1) and (S2). By the rearrangement inequality $U$ is an upper bound
for $L$ on the pair. For the pair formed by the classes $(0,11,6)$ for White and $(0,13,8)$ for
Black, Lemma~\ref{lemplace} and a run over the $3612$ positions of the kings give the exact
maximum, attained at $B_{\ast}$, with $L(B_{\ast})=1250.2997\ldots$. For every other pair we found
$U\leq 1245.67$.
\end{proof}

The values $a=2$ and $b=4$ satisfy~\eqref{eqfeas} for $B_{\ast}$: two pawns and four rooks have
been captured and fourteen pawns promoted. But $B_{\ast}$ is not reachable, for the same reason
as $B_{\max}$: its white king is attacked by four queens. So we must also take checks into account,
and the next lemma states the three facts we use. We say that a man gives a
contact check if it attacks the enemy king from an adjacent square, or if it is a knight
attacking the king.

\begin{lemma}\label{lemcheck}
Consider a position that occurs in a game, after at least one move.
\begin{itemize}
\item[(K1)] The king of the side that is not to move is not attacked.
\item[(K2)] The king of the side to move is attacked by at most two men, and at most one of them
gives a contact check.
\item[(K3)] If a man on a square $t$ gives a contact check, then one of the following squares is
empty: for a rook, bishop, or queen, a square adjacent to $t$ along one of its lines of movement;
for a knight, a square at a knight's move from $t$; for a pawn, and also for an officer standing
on the rank where its pawns are promoted, one of the three squares from which a pawn could have
arrived at $t$ in one step, straight or by capturing.
\end{itemize}
\end{lemma}

\begin{proof}
A player may not leave his own king attacked; this is (K1). For the same reason the king of the
side now to move was not attacked before the last move. So every man attacking it now either
made that move, the new man created by a promotion and the rook in castling included, or is a
rook, bishop, or queen whose line to the king was opened by the move. A line is opened by
vacating a square strictly between the two men, so a man of the second kind is not adjacent to
the king, and a contact check is always given by the man that moved. This proves the second
half of (K2), the king giving no check in castling.

We turn to the first half. A move vacates one square, which lies on at most one line through
the king, so that at most one line is opened. There are two exceptions. Castling vacates two
squares but opens no line: both lie on the back rank, and on that rank the castled king and rook
stand between them. An en passant capture vacates two
squares, that of the capturing pawn and that of the captured one, and may open two lines; but
if the capturing pawn then gives check itself, the king stands diagonally in front of it, two
ranks away from the captured pawn and on an adjacent file, and no line joins these two squares.
In all cases at most two men attack the king.

For (K3), the man giving the contact check has just moved, and the square it left is empty. It
is not a rook that has just castled: the squares next to such a rook along its lines are the
square of the castled king, the square that king has just left, and a square next to the castled
king, and the enemy king can stand on none of them.
For a rook, bishop, or queen, the square next to $t$ in the direction from which it came is the
square it left or a square it passed over. A knight left a square at a knight's move. A pawn came
from the adjacent rank or passed over it, in its own file, or it came from an adjacent file by
capturing; the same holds for the pawn that has just been promoted on $t$.
\end{proof}

\begin{theorem}\label{thmreach}
Let $B_{\circ}$ be the placement with $26$ men
\begin{displaymath}
\begin{array}{ll}
\sq{k}\colon \sq{h8}, & \sq{K}\colon \sq{a8},\\
\sq{q}\colon \sq{g8},\sq{f8},\sq{e8},\sq{d8},\sq{c8},\sq{h7},\sq{g7},\sq{f7},\sq{e7}, &
\sq{Q}\colon \sq{b8},\sq{d7},\sq{c7},\sq{b7},\sq{a7},\sq{h6},\sq{g6},\\
\sq{b}\colon \sq{f6},\sq{e6}, & \sq{B}\colon \sq{d6},\sq{c6},\\
\sq{n}\colon \sq{b6},\sq{a6}, & \sq{N}\colon \sq{h5},\sq{g5} .
\end{array}
\end{displaymath}
Then $B_{\circ}$ is reachable, and $S(B)<S(B_{\circ})$ for every other $B\in\rea$. The code
$S(B_{\circ})$ has $543$ digits, with
\begin{displaymath}
\log_{10}S(B_{\circ})=542.91021\ldots,\qquad L(B_{\circ})=1250.09697\ldots .
\end{displaymath}
In particular $S(B)<10^{543}$ for every $B\in\rea$, and $S_{\mathrm{tot}}<10^{551}$ for every
position occurring in a game.
\end{theorem}

\begin{proof}
We begin with the upper bound. Describe a placement by the $832$ variables $x_{s,l}\in\{0,1\}$, with
$x_{s,l}=1$ if square $s$ carries the label $l$. Conditions (S1) to (S3) and~\eqref{eqfeas}
are linear in these variables, in the excesses, and in $a$. So are (K1) to (K3) once the side to
move is fixed by a further variable: a queen on $t$ attacks a king on $k$ along a line with $j$
squares between them exactly when the $j+2$ corresponding variables are all equal to $1$, which
is expressed by one linear inequality and an auxiliary variable for each triple of a color and two
squares; knights and pawns are simpler. Since $L$ is linear as well, maximizing $L$ under these
necessary conditions for reachability is an integer linear program. We solved it with two
independent programs. The solver HiGHS~\cite{highs}, working in floating point, returns
$B_{\circ}$ with White to move. The solver CP-SAT~\cite{cpsat}, working in integer arithmetic with
the coefficients $\ln p_i$ rounded to nine decimal places, returns $B_{\circ}$ as well, and when
$B_{\circ}$ is excluded its optimum drops by more than $8\cdot 10^{-3}$. The rounding changes $L$
by less than $1.3\cdot 10^{-7}$ on any element of $\adm$, the sum of the labels being at most
$260$. Hence $L(B)<L(B_{\circ})$ for every reachable $B\neq B_{\circ}$.

It remains to show that $B_{\circ}$ is reachable. A game of $272$ plies, with $6$ captures and
$14$ promotions, ends in
$B_{\circ}$ with White to move and in check from the knight on $\sq{b6}$. It was found in two
stages by beam searches over legal moves, the first one leading to fourteen promotions and the
second one rearranging the men without captures. The game is included with the
programs~\cite{software}, and each of its moves was checked for legality with the python-chess
library~\cite{pychess}.

The last assertion follows from $542.911+7.809<551$.
\end{proof}

Compared with Proposition~\ref{propcap}, the check conditions lower $\log_{10}S$ by less than
$0.09$. The material is still that of $B_{\ast}$. What changes is the place of the white king,
which moves from $\sq{f8}$ to the corner $\sq{a8}$, behind three of its own queens.

The same integer program, with the number of empty squares fixed, bounds the code of a
reachable placement with a given number of men.

\begin{proposition}\label{propmen}
For $2\leq m\leq 32$, let $U(m)$ be the largest code of a placement of $\adm$ with exactly $m$
men that satisfies~\eqref{eqfeas} for some integer $a$ and (K1) to (K3) for some side to move.
Then every reachable placement with $m$ men has a code at most $U(m)$, and the values of
$\log_{10}U(m)$ are those of Table~\ref{tabmen}. For $m=2$, $3$, and $26$ the bound is attained
by a reachable placement.
\end{proposition}

\begin{table}[H]
\begin{center}
\begin{tabular}{r|r||r|r||r|r}
$m$ & $\log_{10}U(m)$ & $m$ & $\log_{10}U(m)$ & $m$ & $\log_{10}U(m)$\\ \hline
$2$ & $57.04866$ & $13$ & $319.07923$ & $24$ & $534.27684$\\
$3$ & $81.92005$ & $14$ & $340.38544$ & $25$ & $540.87677$\\
$4$ & $106.42283$ & $15$ & $361.62396$ & $26$ & $542.91021$\\
$5$ & $130.90681$ & $16$ & $382.77724$ & $27$ & $540.38361$\\
$6$ & $155.32538$ & $17$ & $403.91198$ & $28$ & $533.33068$\\
$7$ & $179.65357$ & $18$ & $424.83052$ & $29$ & $480.84857$\\
$8$ & $203.95110$ & $19$ & $445.52020$ & $30$ & $427.25283$\\
$9$ & $228.13943$ & $20$ & $466.17040$ & $31$ & $372.63932$\\
$10$ & $252.23876$ & $21$ & $484.45485$ & $32$ & $317.31944$\\
$11$ & $276.23550$ & $22$ & $502.70312$ &  & \\
$12$ & $297.67365$ & $23$ & $518.50687$ &  & \\
\end{tabular}
\end{center}
\caption{Upper bounds for the code of a reachable placement with $m$ men, rounded down to
five decimal places.}
\label{tabmen}
\end{table}

\begin{proof}
The first assertion follows from Corollary~\ref{corcap} and Lemma~\ref{lemcheck}. We solved
the integer program of Theorem~\ref{thmreach}, with the added condition that exactly $64-m$
squares are empty, for each $m$, with HiGHS and with CP-SAT; the two programs return the
same placement every time. For $m=32$ CP-SAT did not prove optimality within our time limit, so
that value rests on HiGHS alone. For $m=26$ the bound is attained by Theorem~\ref{thmreach}. For
$m=2$ the optimal placement has the white king on $\sq{f8}$ and the black king on $\sq{h8}$, and
for $m=3$ a black queen stands between them on $\sq{g8}$; both are reached by games of about
$50$ plies, found and checked as in Proposition~\ref{propreach} below.
\end{proof}

The bound grows with $m$ up to $m=26$ and then falls, because a full board leaves no room
for promotions. Up to $m=11$ the optimal placement has a lone white king facing a black king
and black queens on the two top ranks; white queens enter from $m=12$ on, and pawns from
$m=29$ on.

\begin{remark}
Bounds of about $307$ and about $800$ digits for the code, and the bound $L<700$, have
appeared in earlier versions of this work that were circulated informally. They are superseded
by Theorems~\ref{thmmax} and~\ref{thmreach}: the first and the third are too small even for
reachable placements, and the second is too large.
\end{remark}

\section{The genome}\label{secgenome}

\begin{definition}
Let $N$ be the number of elements of $\adm$. The chess genome is the increasing
sequence $(a(n))_{1\leq n\leq N}$ of the codes $S(B)$, $B\in\adm$. We write $B_n$ for the
placement with code $a(n)$, and $A(k)$ for the number of $n$ with $a(n)\leq 10^k$.
\end{definition}

By Theorem~\ref{thmorder} the placements $B_1,B_2,\ldots,B_N$ are also those of $\adm$ in
increasing order of $L$. The codes of reachable placements form a subsequence. Its largest term is
given by Theorem~\ref{thmreach}, and Proposition~\ref{propreach} below determines its terms up
to $10^{15}$.

\subsection{First terms and the counting function}

An element of $\adm$ with code at most $X$ is obtained from a pair of kings by adding men one
at a time on squares of increasing index, and every intermediate placement is again in $\adm$
with a smaller code. A depth-first search with exact integer arithmetic therefore lists all codes
up to $X$ in time proportional to their number. Table~\ref{tabfirst} shows the first sixteen terms
and Table~\ref{tabcount} the values of $A(k)$.

\begin{table}[H]
\begin{center}
\begin{tabular}{r|r|l|c}
$n$ & $a(n)$ & $B_n$ & $\log_{10}a(n)$\\ \hline
@@TABFIRST@@
\end{tabular}
\end{center}
\caption{The first sixteen terms of the genome.}
\label{tabfirst}
\end{table}

The genome does not proceed by number of men. The second term is the first one with the kings
exchanged, a white rook on $\sq{b1}$ comes before the white pawn on $\sq{b2}$ because
$3^3<29$, and the placement of two kings on $\sq{a1}$ and $\sq{d1}$ occurs only at $n=8$,
after five placements with three men.

\begin{table}[H]
\begin{center}
\begin{tabular}{r|r||r|r||r|r}
$k$ & $A(k)$ & $k$ & $A(k)$ & $k$ & $A(k)$\\ \hline
@@TABCOUNT@@
\end{tabular}
\end{center}
\caption{The number $A(k)$ of terms of the genome not exceeding $10^k$.}
\label{tabcount}
\end{table}

\subsection{The number of terms}

The number of terms can be computed exactly. We count the placements by choosing the two
kings first, then the pawns, which are confined to ranks $2$ to $7$, and last the officers, which
may stand on any square that is left. The three choices interact only through condition (S3),
which ties the pawns of a color to its officers.

For one color, call a vector $(P,R,N,B,Q)$ of non-negative integers admissible if it
satisfies~\eqref{eqexcess}, and for integers $P,O\geq 0$ put
\begin{displaymath}
g(P,O)=\sum \frac{1}{R!\,N!\,B!\,Q!},
\end{displaymath}
the sum being over the admissible vectors $(P,R,N,B,Q)$ with $R+N+B+Q=O$. Let $K_\kappa$
be the number of ordered pairs of non-adjacent squares of which exactly $\kappa$ lie on ranks
$2$ to $7$; a direct count gives $K_0=212$, $K_1=1448$, and $K_2=1952$.

\begin{theorem}\label{thmcount}
The number $c(m)$ of statically admissible placements with exactly $m$ men is
\begin{equation}\label{eqcount}
c(m)=\sum_{\kappa=0}^{2}K_\kappa\sum
\binom{48-\kappa}{P}\binom{P}{P_1}\frac{(62-P)!}{(62-P-O)!}\,g(P_1,O_1)\,g(P_2,O_2),
\end{equation}
the inner sum being over the non-negative integers $P_1,O_1,P_2,O_2$ with
$P_1+O_1+P_2+O_2=m-2$, and with $P=P_1+P_2$ and $O=O_1+O_2$. The values are given in
Table~\ref{tablayers}, and
\begin{displaymath}
N=\sum_{m=2}^{32}c(m)=@@NTOTAL@@\approx 2.148\cdot 10^{49}.
\end{displaymath}
\end{theorem}

\begin{proof}
Choose the squares of the black and the white king; this can be done in $K_\kappa$ ways with
$\kappa$ kings on ranks $2$ to $7$, and it leaves $62$ free squares, $48-\kappa$ of them
available to pawns. Choose the $P$ squares of the pawns and then the $P_1$ white ones among
them. For given numbers $n_1,\ldots,n_8$ of officers of the eight kinds, with sum $O$, the
officers can be placed on the remaining $62-P$ squares in
$(62-P)!/\bigl((62-P-O)!\,n_1!\cdots n_8!\bigr)$ ways. Summing over the admissible vectors
of the two colors gives~\eqref{eqcount}. The evaluation was done in exact rational arithmetic.
\end{proof}

\begin{table}[H]
\begin{center}
\begin{tabular}{r|r}
$m$ & $c(m)$\\ \hline
@@TABLAYERS@@
\end{tabular}
\end{center}
\caption{The number $c(m)$ of statically admissible placements with $m$ men.}
\label{tablayers}
\end{table}

The first value is $c(2)=K_0+K_1+K_2=3612$, the number of ways to place a white and a black
king on non-adjacent squares. On an $n\times n$ board this number is
$(n-1)(n-2)(n^2+3n-2)$, sequence \seqnum{A035286}. For comparison with $N$, Tromp's
estimate~\cite{tromp} of the number of legal positions, which also distinguishes the side to
move, castling rights, and en passant, is about $4.8\cdot 10^{44}$. Static admissibility is therefore
a generous relaxation of legality.

The next value also has a closed form on an $n\times n$ board with $n\geq 3$, if we keep the
labels and conditions (S1) and (S2), with pawns allowed on ranks $2$ to $n-1$. Write
$c_n(m)$ for the number of such placements with $m$ men, so that $c(m)=c_8(m)$ for $m\leq 3$
and $c_n(2)=(n-1)(n-2)(n^2+3n-2)$.

\begin{proposition}\label{propnn}
For $n\geq 3$,
\begin{displaymath}
c_n(3)=10n^6-4n^5-110n^4+164n^3+92n^2-272n+112 .
\end{displaymath}
The first values, for $n=3,4,\ldots,10$, are $1960$, $19696$, $96552$, $330112$, $904120$,
$2127600$, $4481176$, and $8670592$.
\end{proposition}

\begin{proof}
A placement with three men consists of the two kings and one more man, and condition (S3)
holds trivially. Fix the squares of the kings, and let $\kappa$ be the number of them on ranks $2$
to $n-1$. The third man is one of the eight officers on any of the $n^2-2$ free squares, or one of
the two pawns on any of the $n(n-2)-\kappa$ free squares of ranks $2$ to $n-1$. Summing over the
$c_n(2)$ pairs of squares gives
\begin{displaymath}
c_n(3)=\bigl(8(n^2-2)+2n(n-2)\bigr)\,c_n(2)-2\sum\kappa ,
\end{displaymath}
the sum being over the pairs. A square $s$ on ranks $2$ to $n-1$ is the square of a given king in
exactly $n^2-1-d(s)$ pairs, with $d(s)$ the number of squares adjacent to $s$. On each of these
$n-2$ ranks, the two squares on the edge files have $d(s)=5$ and the other $n-2$ have $d(s)=8$.
Counting both kings,
\begin{displaymath}
\sum\kappa=2\Bigl(n(n-2)(n^2-1)-(n-2)\bigl(10+8(n-2)\bigr)\Bigr)=2(n-2)(n^3-9n+6).
\end{displaymath}
Hence $c_n(3)=(10n^2-4n-16)(n-1)(n-2)(n^2+3n-2)-4(n-2)(n^3-9n+6)$, which expands to
the stated polynomial. For $n=8$ it gives $2127600$, in agreement with Table~\ref{tablayers}.
\end{proof}

\subsection{Reachable terms}

Far from the top, where the placements are sparse, reachability can be decided term by term:
either a short game reaches the placement, or Lemma~\ref{lemcheck} rules it out. The first
term that fails is $a(113)$.

\begin{proposition}\label{propreach}
Of the $433$ terms of the genome not exceeding $10^{15}$, exactly two are not codes of reachable
placements:
\begin{displaymath}
\begin{array}{ll}
a(113)=156600000000000 & (\sq{ka1},\sq{Rb1},\sq{Kc1},\sq{Pb2}),\\
a(178)=391500000000000 & (\sq{Ka1},\sq{Rb1},\sq{kc1},\sq{Pb2}).
\end{array}
\end{displaymath}
In particular, all terms of Table~\ref{tabfirst} are codes of reachable placements.
\end{proposition}

\begin{proof}
In $B_{113}$ the white rook on $\sq{b1}$ attacks the black king on $\sq{a1}$. With White to move
this violates (K1). With Black to move the rook gives a contact check, and the three squares
next to $\sq{b1}$ along the lines of a rook, namely $\sq{a1}$, $\sq{c1}$, and $\sq{b2}$, are occupied;
as White promotes on the eighth rank, this violates (K3). The placement $B_{178}$ is the same
with the kings exchanged, and the same argument applies with the colors of the kings reversed.

For each of the other $431$ terms we found a game that ends in the placement. The terms up to
$10^{15}$ consist of the two kings and at most two other men, so the games mostly consist of
captures. We found them by a beam search over legal moves that keeps the men needed for the
target and steers the others toward capture. When the target contains two pawns on one file, or
a man giving a contact check, the search starts instead from a predecessor obtained by taking
back the last move. For example, $a(205)$, with a black rook on $\sq{b1}$ checking the white king
on $\sq{c1}$, is reached by the move a2xb1=R. The games have at most $106$ plies, are included
with the programs~\cite{software}, and were replayed with python-chess.
\end{proof}

A game for $B_1$ must end with a capture that leaves the two kings on $\sq{a1}$ and $\sq{c1}$,
because a position with two bare kings ends the game at once.

\subsection{Gaps}

One reason to use $L$ is that it is a small real number, so that a database of positions can be
sorted on a floating-point key instead of on integers with hundreds of digits. This works only
if the key separates the placements. If two values of $L$ agree to the working precision, the
two placements are merged or sorted in the wrong order. We therefore need to know how close two
values can be. Small gaps appear early in the genome. The consecutive terms
\begin{displaymath}
\begin{array}{ll}
91815791091791872 & (\sq{Ka1},\sq{kd1},\sq{Pe2},\sq{Pf3}),\\
91815800000000000 & (\sq{ka1},\sq{Kc1},\sq{Pc2},\sq{Pa3},\sq{Pf7})
\end{array}
\end{displaymath}
differ by $9.7\cdot 10^{-8}$ in $L$, the smallest gap among the $10282$ terms below $10^{17}$.
Over the whole of $\adm$ the gaps become far smaller.

\begin{proposition}\label{propgap}
Let $\delta$ be the smallest value of $|L(A)-L(B)|$ over the pairs of distinct elements of
$\adm$. Then
\begin{displaymath}
10^{-138}<\delta<1.91\cdot 10^{-75}.
\end{displaymath}
\end{proposition}

\begin{proof}
A first upper bound comes from counting: the $N$ values of $L$ on $\adm$ are distinct and lie in
the interval from $L(B_1)$ to $L(B_{\max})$, of length less than $1386.43$ by
Corollary~\ref{corbounds}, so two consecutive ones are at distance at most
$1386.43/(N-1)<6.5\cdot 10^{-47}$. An explicit pair does better. The placements
\begin{displaymath}
\begin{array}{l}
\texttt{1KB2NrR/1P2r3/1pqp1rR1/B3P1Q1/1P1r2B1/2k1N1Pp/1pbQ2R1/Nqbn2bn},\\
\texttt{1kb4n/R1qP1RPr/p1n1P2N/1PrN3P/B1n1B1KP/qP1r4/p1Q1R3/3brr2},
\end{array}
\end{displaymath}
with $32$ and $31$ men, are statically admissible, their codes have $302$ digits each, and their
values of $L$ differ by $1.9018\ldots\cdot 10^{-75}$. We found them by looking for short vectors in the
lattice described next, for several values of $C$ between $10^{50}$ and $10^{84}$, and by solving a small integer program that writes
such a vector as the difference of the label vectors of two elements of $\adm$.

For the lower bound, let $A\neq B$ be in $\adm$ with labels $v_i$ and $v_i'$, and put
$c_i=v_i-v_i'$. The replacements used in the proof of Lemma~\ref{lemmat} increase labels, so
$\sum_i v_i$ and $\sum_i v_i^2$ are largest on $\adm$ for the material of $B_{\max}$, with the
values $260$ and $2292$. As $v_iv_i'\geq 0$, this gives
\begin{displaymath}
\sum_i c_i^2\leq 4584\qquad\text{and}\qquad \sum_i |c_i|\leq 520 .
\end{displaymath}
Take $M=256$ and $C=10^{141}$, let $w_i$ be the integer nearest to $C\ln p_i$, and let
$\Lambda\subset\ZZ^{65}$ be the lattice spanned by the vectors $(Me_i,w_i)$ for $1\leq i\leq 64$,
with $e_i$ the unit vectors of $\ZZ^{64}$. The vector $u=\sum_i c_i\,(Me_i,w_i)$ is a nonzero
element of $\Lambda$. Its last coordinate differs from $C\bigl(L(A)-L(B)\bigr)$ by at most
$\frac12\sum_i|c_i|\leq 260$. If we had $|L(A)-L(B)|\leq 1983/C$, the square of the norm of $u$
would be at most $4584M^2+(1983+260)^2$. On the other hand, let $b_1,\ldots,b_{64}$ be any basis of
$\Lambda$ and $b_1^{\ast},\ldots,b_{64}^{\ast}$ its Gram--Schmidt orthogonalization. A nonzero
vector $\sum_i a_ib_i$ of $\Lambda$, with $a_k$ its last nonzero coefficient, has the component
$a_kb_k^{\ast}$ along $b_k^{\ast}$ and is therefore at least as long as $b_k^{\ast}$. We reduced the
basis above with the LLL algorithm~\cite{lll} and its blockwise variant BKZ, as implemented in
fplll~\cite{fplll}, verified in integer arithmetic that the result is again a basis
of $\Lambda$, and computed the squared Gram--Schmidt norms exactly, as quotients of consecutive
leading principal minors of the Gram matrix, with FLINT~\cite{flint}. The smallest of them
exceeds $4584M^2+2243^2$. Hence $\delta>1983\cdot 10^{-141}>10^{-138}$.
\end{proof}

The elementary argument, $S(A)\geq S(B)+1$ for codes in increasing order, only gives
$\delta>1/S_{\max}>10^{-614}$, so the lattice computation gains $476$ orders of magnitude. The
true order of $\delta$ is probably well below the explicit pair: there are about
$2.3\cdot 10^{98}$ pairs of placements, and if their differences were spread evenly over an
interval of length $1386$, the smallest would be about $6\cdot 10^{-96}$. This is only a heuristic.

For computation this means that approximations of $L$ with an absolute error
as small as $10^{-75}$ do not always determine the order of two placements, whereas an
error below $5\cdot 10^{-139}$ always does. Double precision, or even seventy-four decimal places, is
therefore not a safe sorting key for the whole of $\adm$. It is safe to compare floating-point
values of $L$ that come with a rigorous error bound, and to compare the integers $S$ themselves
when the two intervals overlap. These integers have at most $614$ digits, so the exact
comparison costs little.

\section{Comparison with other encodings}\label{secother}

Any injective map from $\adm$ to the positive integers takes a value of at least
$N>2\cdot 10^{49}$, so the code $S$, with values up to $10^{614}$, uses far more digits than necessary. We
compare it with two standard encodings.

The first is the base-$13$ numeral $D(B)=\sum_{i=1}^{64}v_i\,13^{i-1}$, which is a bijection from
all placements onto the integers from $0$ to $13^{64}-1$, and $13^{64}<1.97\cdot 10^{71}$.

The second uses the pairing function of Cantor~\cite{cantor},
\begin{displaymath}
\pi(x,y)=\frac{(x+y)(x+y+1)}{2}+y,
\end{displaymath}
a bijection from pairs of non-negative integers onto the non-negative integers. Starting from
$c_0(i)=v_i$ for $1\leq i\leq 64$, put $c_{r+1}(i)=\pi\bigl(c_r(2i-1),c_r(2i)\bigr)$; six rounds
reduce the $64$ labels to the single integer $C(B)=c_6(1)$, and $C$ is injective because $\pi$
is. As $\pi$ is increasing in each argument, $C(B)$ is at most the value obtained for the
constant vector with all entries $12$, which is less than $1.47\cdot 10^{89}$. The side to move,
the castling rights, and the en passant file can be appended by one more pairing.

\begin{table}[H]
\begin{center}
\begin{tabular}{l|c|c|c}
placement & $S$ & $C$ & $D$\\ \hline
$B_1=(\sq{ka1},\sq{Kc1})$ & $2.0\cdot 10^{11}$ & $6.7\cdot 10^{59}$ & $1871$\\
initial array & $1.2\cdot 10^{238}$ & $6.8\cdot 10^{72}$ & $6.8\cdot 10^{70}$\\
$B_{\max}$ & $2.6\cdot 10^{613}$ & $3.2\cdot 10^{82}$ & $1.9\cdot 10^{71}$
\end{tabular}
\end{center}
\caption{Three encodings of three placements.}
\label{tabcompare}
\end{table}

Table~\ref{tabcompare} shows that neither $S$ nor $C$ is always the smaller one. Iterated pairing
gives much smaller values than $S$ on a full board. But each round roughly squares its input, so
that even two bare kings produce sixty digits, and on sparse boards $S$ is smaller. The numeral
$D$ is smaller than both in the three examples. The advantage of $S$ lies elsewhere: the contents
of a square are a $p$-adic valuation of the code, and a move multiplies the code by a rational
number involving at most four primes. Neither $C$ nor $D$ behaves so simply under moves. For
none of the three encodings does the numerical order have an evident meaning for the game.

\section{Concluding remarks}\label{secopen}

The prime code writes a chess position as an integer in a way that keeps the board visible:
each square can be read off the code, and a move multiplies the code by a small rational
number. Listing the codes in order gives a finite sequence whose basic invariants can be
computed exactly on the statically admissible placements, namely its length, its extreme terms,
its first terms, and the number of terms with a given number of men. For the placements that
occur in games the questions are harder, and we could answer some of them only by combining
short proofs with large searches.

The question we find most interesting concerns the reachable terms. The genome and its
reachable subsequence agree up to $a(112)$, and they still agree on all but two of the first $433$
terms, the two exceptions being placements in which a contact check could not have been given
by any last move. How fast does the reachable subsequence thin out after that? Among the
$2548$ terms up to $10^{16}$, Lemma~\ref{lemcheck} rules out $15$, and for the others we have
not settled the question. At the other end, Theorem~\ref{thmreach} shows that the largest
reachable code has $71$ fewer digits than the largest admissible one. A description of the
reachable terms between these two extremes, even an asymptotic one, would measure how much
the rules of play add to the three static conditions.

An implementation of the encoding, the programs behind the computations of this paper, and all
games quoted in it are available from the repository~\cite{software}.

\section{Acknowledgments}

This paper is dedicated to the memory of Prof.\ Otto Kegel (1934--2025), who was the second
author's tutor at the University of Freiburg.

\begin{thebibliography}{99}

\bibitem{cantor}
G. Cantor, Ein Beitrag zur Mannigfaltigkeitslehre, \emph{J. Reine Angew. Math.} \textbf{84}
(1878), 242--258.

\bibitem{software}
J. A. C\'espedes Bejerano and D. Luebbert, ChessGenome, software repository, 2026. Available at
\url{https://github.com/Yulces/ChessGenome}.

\bibitem{fide}
FIDE, \emph{FIDE Handbook, E. Miscellaneous, 01. Laws of Chess}, 2023. Available at
\url{https://handbook.fide.com}.

\bibitem{pychess}
N. Fiekas, python-chess, software library. Available at
\url{https://github.com/niklasf/python-chess}.

\bibitem{flint}
The FLINT team, FLINT: Fast Library for Number Theory, software library. Available at
\url{https://flintlib.org}.

\bibitem{fplll}
The FPLLL development team, fplll, a lattice reduction library, software library. Available at
\url{https://github.com/fplll/fplll}.

\bibitem{godel}
K. G\"odel, \"Uber formal unentscheidbare S\"atze der Principia Mathematica und verwandter
Systeme I, \emph{Monatsh. Math. Phys.} \textbf{38} (1931), 173--198.

\bibitem{gourion}
D. Gourion, An upper bound for the number of chess diagrams without promotion, \emph{ICGA J.}
\textbf{44} (2022), 44--55.

\bibitem{hardy}
G. H. Hardy and E. M. Wright, \emph{An Introduction to the Theory of Numbers}, 6th edition,
Oxford University Press, 2008.

\bibitem{highs}
Q. Huangfu and J. A. J. Hall, Parallelizing the dual revised simplex method, \emph{Math.
Program. Comput.} \textbf{10} (2018), 119--142.

\bibitem{lll}
A. K. Lenstra, H. W. Lenstra, Jr., and L. Lov\'asz, Factoring polynomials with rational
coefficients, \emph{Math. Ann.} \textbf{261} (1982), 515--534.

\bibitem{oeis}
OEIS Foundation Inc., The On-Line Encyclopedia of Integer Sequences, 2026. Available at
\url{https://oeis.org}.

\bibitem{cpsat}
L. Perron and F. Didier, CP-SAT, software, Google OR-Tools. Available at
\url{https://developers.google.com/optimization/cp/cp_solver}.

\bibitem{shannon}
C. E. Shannon, Programming a computer for playing chess, \emph{Philos. Mag. (7)} \textbf{41}
(1950), 256--275.

\bibitem{steinerberger}
S. Steinerberger, On the number of positions in chess without promotion, \emph{Internat. J.
Game Theory} \textbf{44} (2015), 761--767.

\bibitem{tromp}
J. Tromp, ChessPositionRanking, software repository, 2021. Available at
\url{https://github.com/tromp/ChessPositionRanking}.

\end{thebibliography}

\bigskip
\hrule
\bigskip

\noindent 2020 {\it Mathematics Subject Classification}: Primary 05A15; Secondary 00A08,
11A51, 11Y55, 91A46.

\noindent \emph{Keywords:} chess, prime factorization, G\"odel numbering, enumeration, pairing
function, promotion, integer programming, lattice reduction.

\bigskip
\hrule
\bigskip

\noindent (Concerned with sequences
\seqnum{A000040} and
\seqnum{A035286}.)

\bigskip
\hrule
\bigskip

\end{document}
